\section{Functions}

A function assigns an output to an input. If \(f\) is a function and \(x\) is an input, then
\[
f(x)
\]
denotes the output corresponding to \(x\).

\begin{definitionbox}
A function describes dependence. When we write
\[
y=f(x),
\]
we say that the value of \(y\) is determined by the value of \(x\).
\end{definitionbox}

For example, let
\[
f(x)=x^2+1.
\]
Then
\[
f(0)=1, \qquad f(1)=2, \qquad f(3)=10.
\]

\begin{remarkbox}
The notation \(f(x)\) should not be read as multiplication of \(f\) and \(x\). It means: the value of the function \(f\) at the argument \(x\).
\end{remarkbox}

A function may also depend on more than one variable. For example,
\[
g(x,y)=x^2+y^2
\]
has two inputs. If \(x=3\) and \(y=4\), then
\[
g(3,4)=3^2+4^2=25.
\]

Functions are important not only because they allow us to compute values, but because they show how one quantity depends on another.

\begin{center}
\begin{tabular}{@{}lll@{}}
\toprule
Function & Type of dependence & Effect of increasing \(x\) \\
\midrule
\(f(x)=2x+1\) & linear & output increases at constant rate \\
\(f(x)=x^2\) & nonlinear & change depends on the current value of \(x\) \\
\bottomrule
\end{tabular}
\end{center}
